\chapter{第二讲\quad 犯罪构成}
\label{lecture:02}
\noindent{\small\hyperref[ref:lecture:02]{查阅本讲扩展规则}}\par

\section{第一节\quad 定罪标准：改良版四要件犯罪构成体系}

\paragraph{教材模型\quad 从行为评价到个人责任}
固定行为人 \(a\)、事件 \(x\)、罪名 \(g\) 和待判断的形态 \(s\)，下文省略自变量 \(a,x\)。令 \(\mathrm{Obj}_g\) 表示与保护法益有联系，\(\mathrm{Sub}_g\) 表示行为主体合格。令 \(\mathrm{Act}_{g,s}\)、\(\mathrm{Mens}_{g,s}\) 表示符合该形态的客观、主观要件。\(\mathrm{Just}\) 表示存在正当化事由，\(\mathrm{Resp}_g\) 表示具备责任条件且无排除责任事由。由此定义三个依次收紧的结论。
\begin{align*}
 \mathrm{Prima}_{g,s}
   &\;:=\;\mathrm{Obj}_g\land\mathrm{Sub}_g
                \land\mathrm{Act}_{g,s}\land\mathrm{Mens}_{g,s},\\
 \mathrm{Wrong}_{g,s}
   &\;:=\;\mathrm{Prima}_{g,s}\land\neg\mathrm{Just},\\
 \mathrm{Liab}_{g,s}
   &\;:=\;\mathrm{Wrong}_{g,s}\land\mathrm{Resp}_g .
\end{align*}
\(\mathrm{Prima}\) 是初步符合构成的结论，\(\mathrm{Wrong}\) 是通过正当化检验后的违法结论，\(\mathrm{Liab}\) 是实体上应负刑事责任的结论。这是教材的分层分析，刑法条文通常分别规定要件与例外。责任年龄、责任能力见第17、18条，正当防卫、紧急避险见第20、21条。违法性认识可能性、期待可能性等是教材的归责分析，具体适用有学说边界。

由合取的消去律可证
\[
 \mathrm{Liab}_{g,s}\ \Rightarrow\
 \mathrm{Wrong}_{g,s}\ \Rightarrow\ \mathrm{Prima}_{g,s}.
\]
逆命题均可反驳。取初步符合而有正当防卫的事实，使 \(\mathrm{Prima}=1,\ \mathrm{Just}=1\)，则 \(\mathrm{Wrong}=0\)。取已实施不法行为但未达该罪责任年龄的人，使 \(\mathrm{Wrong}=1,\ \mathrm{Resp}=0\)，则 \(\mathrm{Liab}=0\)。两例分别说明阻却违法和阻却责任的位置不同。写作时把「违法事实」「应负刑责」区分开，可避免把不同阶段的「犯罪」混为一个谓词。

\paragraph{推导\quad 未成年实行人与成年帮助人}
改编教材案例。15周岁的甲主动提出抢夺，成年乙明知计划并在现场接应，两人如约取得符合入罪标准的财物。题设已排除其他罪名、正当化与乙的责任障碍。
\begin{enumerate}
\item 甲实行抢夺，故意与客观行为均能归入抢夺的违法事实，推出 \(\mathrm{Wrong}(\text{甲})\)。
\item 第17条所列14至16周岁责任范围不含普通抢夺。已知甲为15周岁，推出 \(\neg\mathrm{Resp}(\text{甲})\)，故 \(\neg\mathrm{Liab}(\text{甲})\)。
\item 教材在共同犯罪中采用分层分析。先以甲的实行违法事实为帮助的依附对象，再代入乙的帮助故意和实际接应，推出乙的帮助违法事实。乙的年龄与能力另审，题设均满足，推出乙可按帮助犯负刑责。甲主动并自主实行，乙接应这一事实也不足以推出乙支配甲。
\end{enumerate}
乙若只是碰巧路过、既不知情也未提供帮助，步骤3即缺少独立的主客观条件。即便甲的违法事实保持不变，乙的结论仍须重新推导。

\paragraph{边界\quad 客观危害先审，犯罪形态另选}
客观要件 \(\mathrm{Act}_{g,s}\) 随形态 \(s\) 变化。杀人既遂需有死亡及相应归责，未遂另考察着手、未得逞及意志以外原因，不能照搬既遂的死亡要件。教材以向荒地中的稻草人射击说明争议。若依教材所采的客观危险立场，且题设任何人的生命均无现实危险，\(\mathrm{Act}_{g,\text{未遂}}\) 的危险前提不成立，杀人意思无法填补它。若题设另采以行为人认识为重的危险判断立场，须把该前提明确写出，再讨论能否认定未遂。先固定观点及风险判据，才能比较结果。

\section{第二节\quad 定罪方法：三段论推理}

\paragraph{三段论的构造}
大前提是确定了法条版本及解释口径的条件规则，小前提由经证明的事实归摄为要件，最后作演绎。令 \(F_i\) 为证据支持的具体事实，\(P_i\) 为法律评价后的第 \(i\) 个必要要素。典型证明写作
\begin{align*}
 F_1,\mathcal R&\vdash P_1,\quad \ldots,\quad F_n,\mathcal R\vdash P_n,\\
 P_1,\ldots,P_n&\vdash P_1\land\cdots\land P_n,\\
 P_1\land\cdots\land P_n,\quad
 (P_1\land\cdots\land P_n)\Rightarrow Q&\vdash Q.
\end{align*}
归摄本身需要理由。例如「视频显示持刀」是事实，「正在实施暴力」是结合时间、动作和对象作出的规范判断。教材把要素分为记述与规范、成文与不成文、积极与消极、客观与主观。几组分类可交叉。解释出的不成文要素仍须有法条和解释依据。一个否定性要素如「未经同意」也应接受证据审查，不能以题干遗漏直接认定。

\paragraph{推导\quad 从证据到结论}
自拟题设。甲无权取得乙置于锁柜的财物，秘密破锁取走并带回住所。录像与后续处置证据均已查证，能证明甲明知财物属乙并意图据为己有。数额达到题设法定门槛，其他责任条件具备，也无正当化事实。按第264条的教材分析，破锁并取走归入窃取，乙原有控制与甲建立新控制归入占有转移，后续处置与明知事实支持非法占有目的。分别得到行为、对象、控制转移、主观目的和门槛各项，再合取推出盗窃既遂，最后审查责任与证据。反向由「财物后来在甲处」推出全部条件无效。同样事实也可能由乙寄存造成，此时「未经许可窃取」这一步存在替代解释。

\paragraph{证明模型\quad 存疑如何进入推理}
《刑事诉讼法》第55条要求定罪量刑事实有证据、证据经法定程序查证属实，并综合排除合理怀疑。第200条第3项规定证据不足时作无罪判决。用非空的 \(\Omega_\Gamma\) 表示从经审查证据 \(\Gamma\) 出发仍无法排除的合理事实版本。在这个练习模型中，
\[
 \mathrm{Established}_{\Gamma}(P)
 \;:=\; \forall\omega\in\Omega_\Gamma,\ P(\omega).
\]
\(\Omega_\Gamma\) 仅包含经证据检验的合理版本，不把任意想象都列入。当 \(\Omega_\Gamma\) 为空，应先排查证据冲突或模型输入错误，避免空集上全称命题为真造成虚假定罪。设普通构成 \(Q\) 与加重条件 \(H\) 可分，则高一档需证立 \(Q\land H\)。若所有合理版本都满足 \(Q\)，其中一版本满足 \(H\)、另一版本不满足，能证明 \(Q\)，不能证明 \(Q\land H\)。若尚有合理版本使 \(Q\) 不成立，则基础罪亦不能证明。证据不足推出的是\emph{证明结论}，不赋予某个事实的反面以真实性。仅列出两个互不包含的轻、重罪名，还不足以按较轻者定罪，须独立验证较轻罪的全套要件。

\paragraph{边界\quad 一个事实可接受多次检验}
在同一事实集上分别运行 \(g_1,g_2\) 的证明，可能都得到符合。由 \(\mathrm{Liab}_{g_1}(x)\land\mathrm{Liab}_{g_2}(x)\) 无法直接推出应对两个罪都单独宣告并罚。下一步须检查行为数、法条关系与特别处断规则。教材所举一行为同时触犯数罪的例型，应把构成判定与竞合处断分开写。

\subsection*{闭卷复述与检验}
\textbf{口诀}\quad「客体主体，客观主观。正当化看事，归责另看人。法定规则逐项套，疑点放回证据链。」

\textbf{复述路线}\quad
先标记罪名与形态，再依次写四要件和两个排除层。为每个要件写出一个可核查的事实来源。举一个合理反事实版本，测试它会击穿基础要件、仅击穿加重情节，还是完全不改变结论。最后才讨论竞合和量刑。

\begin{enumerate}
\item 甲为制止正在进行的严重暴力侵害而依法实施正当防卫，致侵害人受伤。究竟哪一步排除刑责？\\
\textbf{答}\quad 初步构成可有伤害事实，成立正当防卫使 \(\mathrm{Just}=1\)，故 \(\mathrm{Wrong}=0\)。无需以伤害结果否定防卫。
\item 甲只供称自己偷过东西，没有其他证据。能否因供述承认全部要件而判有罪？\\
\textbf{答}\quad 不能。要件推理在假定供述真实时可能有效，事实仍未达到第55条的证明要求。
\item 普通要件在两个合理事实版本下都已证立，加重事实只在一个版本成立。应写什么结论？\\
\textbf{答}\quad 在其余条件满足时可以认定普通构成，无法以现有证据适用加重分支。不能写成「加重事实一定没有发生」。
\end{enumerate}
